% ============================================================
% Übungsblatt 2: Floating-Point Arithmetic (Gleitkommaarithmetik)
% zum Vorlesungsskript "Numerical Methods" (Prof. Mark Schutera),
% Kapitel 2 (Original: S. 17-24). Eigenständig kompilierbar.
% ============================================================
\documentclass[a4paper,11pt]{scrartcl}

\usepackage[T1]{fontenc}
\usepackage[utf8]{inputenc}
\usepackage[ngerman]{babel}
\usepackage{lmodern}
\usepackage{amsmath}
\usepackage{amssymb}
\usepackage[a4paper,margin=2.5cm]{geometry}
\usepackage{enumitem}
\usepackage{booktabs}
\usepackage[hidelinks]{hyperref}
\usepackage{catppuccin-dark}

\title{Übungsblatt 2: Floating-Point Arithmetic (Gleitkommaarithmetik)}
\subtitle{zum Vorlesungsskript \glqq Numerical Methods\grqq{} (Prof.~Mark Schutera), Kapitel~2}
\author{}
\date{\today}

% Aufgaben-Zähler: \ref{auf:N} liefert die Aufgabennummer
\newcounter{aufgabe}
\newcommand{\aufgabe}[2]{%
  \refstepcounter{aufgabe}%
  \subsection*{Aufgabe \theaufgabe\ (#2): #1}%
}

\begin{document}

\maketitle

\noindent
Dieses Übungsblatt deckt Kapitel~2 des Skripts ab: Fehlerarten, Gleitkommadarstellung (floating-point representation) nach IEEE~754, Maschinengenauigkeit (machine epsilon), Festkommadarstellung (fixed-point representation) sowie Auslöschung (loss of significance / catastrophic cancellation). \textbf{Gesamtpunktzahl: 38 Punkte.} Die Aufgaben steigen im Schwierigkeitsgrad an: Aufgaben 1--3 prüfen das Begriffsverständnis, Aufgaben 4--7 sind Standardaufgaben analog zum Skript, Aufgabe~8 ist eine Transferaufgabe. Die Lösungen befinden sich ab Seite~\pageref{sec:loesungen} --- erst nach eigener Bearbeitung lesen!

\section*{Aufgaben}

% ------------------------------------------------------------
\aufgabe{Fehlerarten erkennen und zuordnen}{4 Punkte}\label{auf:1}

\begin{enumerate}[label=(\alph*)]
  \item Definieren Sie kurz die drei im Skript unterschiedenen Fehlerarten: Modellierungsfehler (modeling error), Abschneide-/Trunkierungsfehler (truncation error) und Rundungsfehler (rounding error). \hfill (2~Punkte)
  \item Ordnen Sie die folgenden Szenarien jeweils genau einer der drei Fehlerarten zu und begründen Sie kurz: \hfill (2~Punkte)
  \begin{enumerate}[label=(\roman*)]
    \item Ein schwingendes Pendel wird in einer Simulation als reibungsfrei angenommen, obwohl in der Realität Luftwiderstand wirkt.
    \item Die unendliche Reihe $e^x = \sum_{k=0}^{\infty} \frac{x^k}{k!}$ wird zur Berechnung nach dem zehnten Summanden abgebrochen.
    \item Die Zahl $\frac{1}{3}$ wird im Computer als $0.33333333$ gespeichert.
  \end{enumerate}
\end{enumerate}

% ------------------------------------------------------------
\schreibraum[7]
\aufgabe{Warum $10 \div 3$ nicht exakt darstellbar ist}{4 Punkte}\label{auf:2}

\begin{enumerate}[label=(\alph*)]
  \item Erklären Sie anhand der schriftlichen Division, warum $\frac{10}{3}$ mit endlich vielen Dezimalziffern niemals exakt aufgeschrieben oder gespeichert werden kann. \hfill (1~Punkt)
  \item Berechnen Sie den Rundungsfehler (rounding error), der entsteht, wenn man $\frac{10}{3}$ nach drei Nachkommastellen abschneidet (truncate), also als $3.333$ speichert. Geben Sie den Fehler auch als Bruch an. \hfill (2~Punkte)
  \item Geben Sie an, wie der Wert $\frac{10}{3}$ in den folgenden vier Darstellungstypen gespeichert würde: Integer (int), Floating-point (float, 7~Nachkommastellen), Double precision (double, 16~Nachkommastellen), Fixed-point mit 4~Nachkommastellen. \hfill (1~Punkt)
\end{enumerate}

% ------------------------------------------------------------
\schreibraum[7]
\aufgabe{Relative vs.\ absolute Präzision}{3 Punkte}\label{auf:3}

\begin{enumerate}[label=(\alph*)]
  \item Erklären Sie den Unterschied zwischen der relativen und der absoluten Präzision von Gleitkommazahlen. Welche der beiden ist (näherungsweise) konstant, welche hängt von der Größe der dargestellten Zahl ab? Geben Sie die zugehörigen Ausdrücke aus dem Skript an. \hfill (1.5~Punkte)
  \item Das Skript veranschaulicht den Unterschied mit dem Vergleich $1:2 = 0.5$, aber $10:2 = 5$. Erläutern Sie, was diese Randnotiz über Schritte (relative Präzision) und Lücken (gaps, absolute Präzision) aussagt. \hfill (0.5~Punkte)
  \item Berechnen Sie näherungsweise die absolute Präzision (Lückengröße zwischen benachbarten darstellbaren Zahlen) in IEEE~754 single precision ($\varepsilon_{\text{mach}} \approx 1.19 \times 10^{-7}$) für $x = 1$ und für $x = 10^6$. \hfill (1~Punkt)
\end{enumerate}

% ------------------------------------------------------------
\schreibraum[6]
\aufgabe{IEEE 754 single precision: Dekodieren und Kodieren}{6 Punkte}\label{auf:4}

In IEEE~754 single precision (32~Bit: 1~Vorzeichenbit + 8~Exponentenbits + 23~Mantissenbits) wird eine normalisierte Zahl gespeichert als
\[
x = (-1)^s \cdot 1.m \cdot 2^{E - 127},
\]
wobei $s$ das Vorzeichenbit (sign bit), $m$ das Mantissenfeld (mantissa) und $E$ der gespeicherte Exponent mit Bias~127 ist (die führende~$1$ der normalisierten Mantisse wird nicht mitgespeichert). \emph{Hinweis: Die Konvention des impliziten führenden Mantissenbits geht leicht über die Skriptdarstellung (S.~19--20) hinaus und ist hier vollständig angegeben --- eine kleine Transferkomponente.}

\begin{enumerate}[label=(\alph*)]
  \item \textbf{Dekodieren:} Welche Dezimalzahl ist hier gespeichert? \hfill (3~Punkte)
  \[
  s = 1, \qquad E = 10000010_2, \qquad m = 10100000000000000000000_2
  \]
  \item \textbf{Kodieren:} Bestimmen Sie Vorzeichenbit $s$, Exponentenfeld $E$ (binär, 8~Bit) und Mantissenfeld $m$ (binär, 23~Bit) für die Dezimalzahl $x = 6.25$. \hfill (3~Punkte)
\end{enumerate}

% ------------------------------------------------------------
\schreibraum[9]
\aufgabe{Machine Epsilon und die Mantisse}{5 Punkte}\label{auf:5}

\begin{enumerate}[label=(\alph*)]
  \item Definieren Sie Machine Epsilon ($\varepsilon_{\text{mach}}$, Maschinengenauigkeit) und geben Sie die Formel in Abhängigkeit von der Anzahl $t$ der Mantissenbits an. Berechnen Sie daraus $\varepsilon_{\text{mach}}$ für IEEE~754 single precision ($t = 23$) und double precision ($t = 52$). \hfill (2~Punkte)
  \item Betrachten Sie eine (unnormalisierte) 2-Bit-Mantisse mit Significand $0.xx_2$. Listen Sie alle vier Bitkombinationen mit ihren Dezimalwerten auf und geben Sie das zugehörige Machine Epsilon an. \hfill (2~Punkte)
  \item Wie viele verschiedene Werte zwischen zwei benachbarten Zweierpotenzen erlaubt eine Mantisse mit $t$~Bits allgemein? Wie viele sind es konkret bei $t = 3$? \hfill (1~Punkt)
\end{enumerate}

% ------------------------------------------------------------
\schreibraum[8]
\aufgabe{Constructing scale: Zehnerpotenzen in Binärdarstellung}{4 Punkte}\label{auf:6}

Stellen Sie die Zehnerpotenzen $10^1$, $10^2$ und $10^3$ jeweils
\begin{enumerate}[label=(\alph*)]
  \item als Binärzahl dar, \hfill (1.5~Punkte)
  \item in normalisierter Form $1.xxx_2 \times 2^e$ dar, \hfill (1.5~Punkte)
  \item und geben Sie das zugehörige gespeicherte Exponentenfeld $E$ (binär, 8~Bit, Bias~127, IEEE~754 single precision) an. \hfill (1~Punkt)
\end{enumerate}

% ------------------------------------------------------------
\schreibraum[7]
\aufgabe{Fixed-Point-Skalierung}{4 Punkte}\label{auf:7}

Es sollen reelle Zahlen zwischen $-1000$ und $1000$ mit einem 32-Bit signed Integer (Wertebereich $-2147483648$ bis $2147483647$) in Festkommadarstellung (fixed-point representation) gespeichert werden. Dazu wird jede reelle Zahl mit dem Skalierungsfaktor $10^6$ multipliziert und als Integer abgelegt.

\begin{enumerate}[label=(\alph*)]
  \item Welcher Integer-Wert wird für die reelle Zahl $1.234567$ gespeichert? \hfill (0.5~Punkte)
  \item Wie groß ist die kleinste darstellbare Differenz (Präzision) und wie groß ist der maximale Rundungsfehler beim Speichern? \hfill (1.5~Punkte)
  \item Zeigen Sie rechnerisch, dass der Skalierungsfaktor $10^6$ für den geforderten Wertebereich in einen 32-Bit signed Integer passt, $10^7$ aber nicht mehr. \hfill (1~Punkt)
  \item Worin unterscheidet sich die Präzision der Festkommadarstellung grundsätzlich von der der Gleitkommadarstellung? \hfill (1~Punkt)
\end{enumerate}

% ------------------------------------------------------------
\schreibraum[7]
\aufgabe{Transfer: Auslöschung, Overflow/Underflow und Quantisierung}{8 Punkte}\label{auf:8}

\begin{enumerate}[label=(\alph*)]
  \item Eine Maschine speichert Zahlen mit 8~signifikanten Stellen (kaufmännisch gerundet). Gegeben sind die beiden Paare
  \[
  \text{(i)}\quad a = 12345678.5,\ b = 12345678.0
  \qquad\text{und}\qquad
  \text{(ii)}\quad a = 12345678.4,\ b = 12345678.0.
  \]
  Berechnen Sie für beide Paare: die gespeicherten Werte $\tilde{a}$, $\tilde{b}$, das Maschinenergebnis $\tilde{a} - \tilde{b}$, die wahre Differenz $a - b$ sowie den absoluten und den relativen Fehler des Maschinenergebnisses. Was fällt beim Vergleich der beiden Paare auf? \hfill (3~Punkte)
  \item Betrachten Sie die Funktion $f(\epsilon) = \dfrac{1}{1 + \epsilon - 1}$. Was liefert eine Auswertung in IEEE~754 double precision ($\varepsilon_{\text{mach}} \approx 2.22 \times 10^{-16}$) für $\epsilon = 10^{-20}$, wenn der Nenner exakt so wie angegeben berechnet wird --- und warum? Geben Sie eine mathematisch äquivalente Umformulierung an, die das Problem vollständig vermeidet, und werten Sie diese für $\epsilon = 10^{-20}$ aus. \hfill (3~Punkte)
  \item Die größte darstellbare Zahl in IEEE~754 single precision liegt bei ca.\ $10^{38}$, gesetzt durch den größten Exponenten $e = +127$. Was geschieht in single precision bei der Berechnung von $x^2$ für $x = 10^{20}$ und bei der Berechnung von $y^2$ für $y = 10^{-25}$? Benennen Sie beide Phänomene und beschreiben Sie eine gefährliche Folgerechnung, die sich aus dem zweiten Fall ergeben kann. \hfill (1.5~Punkte)
  \item Quantisierte Modelle (quantized models) auf Edge-Geräten rechnen mit niedrigpräziser Arithmetik (z.\,B.\ 8-Bit-Integer); bei binären neuronalen Netzen (binary neural networks, BNNs) sind Gewichte und Aktivierungen sogar auf $\{-1, +1\}$ beschränkt. Erläutern Sie kurz, warum gerade dort das Verständnis von Gleitkommadarstellung und Rundungseffekten kritisch ist. \hfill (0.5~Punkte)
\end{enumerate}

\bigskip
\noindent\textbf{Summe: 38 Punkte}

% ============================================================
\schreibraum[11]
\newpage
\section*{Lösungen --- erst nach Bearbeitung lesen!}\label{sec:loesungen}

% ------------------------------------------------------------
\subsection*{Lösung zu Aufgabe~\ref{auf:1}}

\begin{enumerate}[label=(\alph*)]
  \item
  \begin{itemize}
    \item \textbf{Modellierungsfehler (modeling error):} entsteht, weil alle Modelle Vereinfachungen der Realität sind; die Differenz zwischen Modell und realem System führt einen Fehler ein --- der feine Unterschied zwischen dem, was wir $f(x)$ \emph{nennen}, und dem, was tatsächlich ein $f(x)$ \emph{ist}.
    \item \textbf{Abschneide-/Trunkierungsfehler (truncation error):} entsteht, wenn ein unendlicher Prozess (Reihe, Grenzwert, kontinuierliche Funktion) durch einen endlichen approximiert wird --- auch Approximationsfehler (approximation error) genannt.
    \item \textbf{Rundungsfehler (rounding error):} Differenz zwischen dem wahren Wert und dem Wert, den man nach dem Abschneiden der (Dezimal-)Entwicklung erhält; entsteht, weil Computer Zahlen stets mit endlich vielen Ziffern bzw.\ Bits speichern.
  \end{itemize}
  \item
  \begin{enumerate}[label=(\roman*)]
    \item \textbf{Modellierungsfehler:} Das Modell (reibungsfreies Pendel) vereinfacht die Realität (mit Luftwiderstand); der Fehler entsteht bereits vor jeder Berechnung in der Modellbildung.
    \item \textbf{Trunkierungsfehler:} Ein unendlicher Prozess (unendliche Reihe) wird durch einen endlichen (10~Summanden) ersetzt.
    \item \textbf{Rundungsfehler:} Der Speicherwert $0.33333333$ schneidet die unendliche Dezimalentwicklung $0.\overline{3}$ nach endlich vielen Ziffern ab; der Fehler entsteht durch die endliche Speicherdarstellung.
  \end{enumerate}
\end{enumerate}

% ------------------------------------------------------------
\subsection*{Lösung zu Aufgabe~\ref{auf:2}}

\begin{enumerate}[label=(\alph*)]
  \item Die schriftliche Division liefert in jedem Schritt dasselbe Muster: $10 \div 3 = 3$, Rest~$1$; nach \glqq bring down a 0\grqq{} wieder $10 \div 3 = 3$, Rest~$1$ --- und so weiter. Der Rest wird nie null, also wiederholt sich die Ziffer~$3$ unendlich oft: $\frac{10}{3} = 3.333\ldots = 3.\overline{3}$. Da man beim Aufschreiben oder Speichern irgendwo stoppen muss, entsteht zwangsläufig ein Fehler --- egal, wie viele Ziffern man schreibt.
  \item
  \[
  \text{Rundungsfehler} = \left| 3.333333\ldots - 3.333 \right| = 0.000333\ldots
  \]
  Als Bruch:
  \[
  \frac{10}{3} - \frac{3333}{1000} = \frac{10000 - 9999}{3000} = \frac{1}{3000} = 0.000\overline{3}.
  \]
  \item
  \begin{center}
  \begin{tabular}{ll}
  \toprule
  Darstellung & gespeicherter Wert \\
  \midrule
  Integer (int) & $3$ \\
  Floating-point (float) & $3.3333333$ \\
  Double precision (double) & $3.3333333333333333$ \\
  Fixed-point (4 Nachkommastellen) & $3.3333$ \\
  \bottomrule
  \end{tabular}
  \end{center}
\end{enumerate}

% ------------------------------------------------------------
\subsection*{Lösung zu Aufgabe~\ref{auf:3}}

\begin{enumerate}[label=(\alph*)]
  \item Die \textbf{relative Präzision} von Gleitkommazahlen ist näherungsweise konstant: Sie beträgt ca.\ $2^{-t}$ (mit $t$ = Anzahl der Mantissenbits), unabhängig von der Größe der dargestellten Zahl. Die \textbf{absolute Präzision} hängt dagegen von der Größenordnung der Zahl ab:
  \[
  \text{absolute Präzision} = \varepsilon_{\text{mach}} \cdot |x|.
  \]
  Der absolute Abstand benachbarter darstellbarer Zahlen wächst also proportional zu $|x|$: Große Zahlen haben größere absolute Lücken (gaps) als kleine Zahlen.
  \item Wird das Intervall $[0,1]$ in zwei Schritte geteilt, beträgt die Schrittweite $0.5$; wird das Intervall $[0,10]$ ebenfalls in zwei Schritte geteilt, beträgt sie $5$. Gleiche \emph{Anzahl} von Schritten (gleiche relative Präzision), aber Lücken von $0.5$ vs.\ $5$ --- die absolute Präzision skaliert mit der Größe der Zahl.
  \item Mit $\varepsilon_{\text{mach}} \approx 1.19 \times 10^{-7}$:
  \begin{align*}
  x = 1:&\quad \varepsilon_{\text{mach}} \cdot |x| \approx 1.19 \times 10^{-7} \cdot 1 = 1.19 \times 10^{-7},\\
  x = 10^6:&\quad \varepsilon_{\text{mach}} \cdot |x| \approx 1.19 \times 10^{-7} \cdot 10^6 = 0.119 \approx 0.12.
  \end{align*}
  Bei $x = 10^6$ liegen benachbarte darstellbare Zahlen also bereits gut eine Zehntel-Einheit auseinander, während die Lücke bei $x = 1$ nur etwa $10^{-7}$ beträgt.
\end{enumerate}

% ------------------------------------------------------------
\subsection*{Lösung zu Aufgabe~\ref{auf:4}}

\begin{enumerate}[label=(\alph*)]
  \item \textbf{Dekodieren.} Gespeicherter Exponent:
  \[
  E = 10000010_2 = 128 + 2 = 130_{10}, \qquad e = E - 127 = 130 - 127 = 3.
  \]
  Significand (mit impliziter führender~$1$): das Mantissenfeld beginnt mit $101$, alle weiteren Bits sind null, also
  \[
  1.101_2 = 1 + 1 \cdot 2^{-1} + 0 \cdot 2^{-2} + 1 \cdot 2^{-3} = 1 + 0.5 + 0.125 = 1.625.
  \]
  Mit Vorzeichenbit $s = 1$:
  \[
  x = (-1)^1 \cdot 1.625 \cdot 2^3 = -1.625 \cdot 8 = \mathbf{-13.0}.
  \]
  \emph{Probe:} $13_{10} = 1101_2 = 1.101_2 \times 2^3$.~\checkmark
  \item \textbf{Kodieren von $x = 6.25$.} Umwandlung ins Binärsystem:
  \[
  6.25 = 4 + 2 + 0.25 = 110.01_2.
  \]
  Normalisieren (Komma um zwei Stellen nach links verschieben):
  \[
  110.01_2 = 1.1001_2 \times 2^2 \quad\Rightarrow\quad e = 2.
  \]
  Damit:
  \begin{align*}
  s &= 0 \quad (\text{Zahl ist positiv}),\\
  E &= e + 127 = 2 + 127 = 129_{10} = 10000001_2,\\
  m &= 10010000000000000000000_2 \quad (\text{Nachkommabits } 1001 \text{, mit Nullen auf 23 Bit aufgefüllt}).
  \end{align*}
  \emph{Probe:} $(-1)^0 \cdot 1.1001_2 \cdot 2^2 = (1 + 0.5 + 0.0625) \cdot 4 = 1.5625 \cdot 4 = 6.25$.~\checkmark
\end{enumerate}

% ------------------------------------------------------------
\subsection*{Lösung zu Aufgabe~\ref{auf:5}}

\begin{enumerate}[label=(\alph*)]
  \item Machine Epsilon ist die kleinste positive Zahl, sodass $1 + \varepsilon_{\text{mach}} \neq 1$ in der Computerarithmetik gilt; es quantifiziert die obere Schranke des relativen Fehlers durch Rundung in der Gleitkomma-Arithmetik. Formel:
  \[
  \varepsilon_{\text{mach}} = 2^{-t},
  \]
  mit $t$ = Anzahl der Mantissenbits. Konkret:
  \begin{align*}
  \text{single precision } (t = 23):&\quad \varepsilon_{\text{mach}} = 2^{-23} = \tfrac{1}{8388608} \approx 1.19 \times 10^{-7},\\
  \text{double precision } (t = 52):&\quad \varepsilon_{\text{mach}} = 2^{-52} \approx 2.22 \times 10^{-16}.
  \end{align*}
  \item Alle vier Bitkombinationen der unnormalisierten 2-Bit-Mantisse $0.xx_2$:
  \begin{align*}
  00:\ & 0.00_2 = 0 + 0 \times 2^{-1} + 0 \times 2^{-2} = 0.0,\\
  01:\ & 0.01_2 = 0 + 0 \times 2^{-1} + 1 \times 2^{-2} = 0.25,\\
  10:\ & 0.10_2 = 0 + 1 \times 2^{-1} + 0 \times 2^{-2} = 0.5,\\
  11:\ & 0.11_2 = 0 + 1 \times 2^{-1} + 1 \times 2^{-2} = 0.75.
  \end{align*}
  Machine Epsilon: $\varepsilon_{\text{mach}} = 2^{-2} = 0.25$ --- genau der Abstand zwischen zwei benachbarten Mantissenwerten.
  \item Eine Mantisse mit $t$~Bits erlaubt $2^t$ verschiedene Werte zwischen zwei benachbarten Zweierpotenzen (skaliert durch den Exponenten). Für $t = 3$: $2^3 = 8$~Werte.
\end{enumerate}

% ------------------------------------------------------------
\subsection*{Lösung zu Aufgabe~\ref{auf:6}}

\begin{enumerate}[label=(\alph*)]
  \item Binärdarstellungen:
  \begin{align*}
  10^1 = 10_{10} &= 1010_2 && (8 + 2 = 10),\\
  10^2 = 100_{10} &= 1100100_2 && (64 + 32 + 4 = 100),\\
  10^3 = 1000_{10} &= 1111101000_2 && (512 + 256 + 128 + 64 + 32 + 8 = 1000).
  \end{align*}
  \item Normalisierte Form (Komma hinter die führende~$1$ schieben):
  \begin{align*}
  10^1 &= 1.010_2 \times 2^3,\\
  10^2 &= 1.100100_2 \times 2^6,\\
  10^3 &= 1.111101000_2 \times 2^9.
  \end{align*}
  \item Gespeicherter Exponent $E = e + 127$:
  \begin{align*}
  e = 3:&\quad E = 130_{10} = 10000010_2 && (128 + 2 = 130),\\
  e = 6:&\quad E = 133_{10} = 10000101_2 && (128 + 4 + 1 = 133),\\
  e = 9:&\quad E = 136_{10} = 10001000_2 && (128 + 8 = 136).
  \end{align*}
\end{enumerate}

% ------------------------------------------------------------
\subsection*{Lösung zu Aufgabe~\ref{auf:7}}

\begin{enumerate}[label=(\alph*)]
  \item Gespeichert wird $1.234567 \cdot 10^6 = 1234567$.
  \item Die kleinste darstellbare Differenz ist ein Integer-Schritt, zurückskaliert mit $10^{-6}$:
  \[
  \text{Präzision} = 10^{-6} = 0.000001.
  \]
  Der maximale Rundungsfehler beim Speichern ist ein halber Schritt:
  \[
  0.5 \times 10^{-6} = 0.0000005.
  \]
  \item Größter zu speichernder Betrag: $1000 \cdot 10^6 = 10^9 = 1\,000\,000\,000 < 2\,147\,483\,647$ --- passt (ebenso $-10^9 > -2\,147\,483\,648$). Mit Skalierungsfaktor $10^7$ wäre dagegen $1000 \cdot 10^7 = 10^{10} = 10\,000\,000\,000 > 2\,147\,483\,647$ --- Überlauf des Integer-Bereichs. $10^6$ ist also die größte Zehnerpotenz, die als Skalierungsfaktor möglich ist.
  \item Bei der Festkommadarstellung ist die Lücke zwischen darstellbaren Zahlen (die absolute Präzision) \emph{immer gleich}, egal wie groß oder klein der Wert ist --- das gibt mehr Kontrolle und vorhersagbare Präzision und Performance. Bei der Gleitkommadarstellung ist dagegen die \emph{relative} Präzision näherungsweise konstant ($\approx 2^{-t}$), während die absolute Lücke proportional zu $|x|$ wächst --- dafür ist der darstellbare Wertebereich (Dynamikumfang) viel größer.
\end{enumerate}

% ------------------------------------------------------------
\subsection*{Lösung zu Aufgabe~\ref{auf:8}}

\begin{enumerate}[label=(\alph*)]
  \item \textbf{Paar (i):} $a = 12345678.5$, $b = 12345678.0$. Mit 8~signifikanten Stellen (kaufmännisch gerundet, $.5$ rundet auf):
  \[
  \tilde{a} = 12345679, \qquad \tilde{b} = 12345678.
  \]
  Maschinenergebnis: $\tilde{a} - \tilde{b} = 1$. Wahre Differenz: $a - b = 0.5$.
  \[
  \text{absoluter Fehler} = |1 - 0.5| = 0.5, \qquad
  \text{relativer Fehler} = \frac{0.5}{0.5} = 1 = 100\,\%.
  \]
  \textbf{Paar (ii):} $a = 12345678.4$, $b = 12345678.0$. Mit 8~signifikanten Stellen ($12345678.4$ rundet \emph{ab}):
  \[
  \tilde{a} = 12345678, \qquad \tilde{b} = 12345678.
  \]
  Maschinenergebnis: $\tilde{a} - \tilde{b} = 0$. Wahre Differenz: $a - b = 0.4$.
  \[
  \text{absoluter Fehler} = |0 - 0.4| = 0.4, \qquad
  \text{relativer Fehler} = \frac{0.4}{0.4} = 1 = 100\,\%.
  \]
  \textbf{Beobachtung:} Die wahren Resultate unterscheiden sich nur um $0.1$ ($0.5$ vs.\ $0.4$), die Maschinenresultate aber um eine ganze Einheit ($1$ vs.\ $0$) --- in beiden Fällen beträgt der relative Fehler $100\,\%$. Beim Subtrahieren nahezu gleicher Zahlen heben sich die führenden (korrekten) Ziffern auf, und übrig bleiben nur die rundungsfehleranfälligen letzten Stellen: Auslöschung (loss of significance / catastrophic cancellation) verstärkt den Rundungsfehler massiv.
  \item Da $\epsilon = 10^{-20} < \varepsilon_{\text{mach}} \approx 2.22 \times 10^{-16}$, wird die Addition $1 + \epsilon$ in double precision exakt auf $1$ gerundet ($1 + \epsilon$ und $1$ sind nicht unterscheidbar). Der Nenner wird damit
  \[
  \underbrace{(1 + \epsilon)}_{\to\, 1} - 1 = 0,
  \]
  und die Auswertung liefert eine Division durch null (in IEEE~754: $+\infty$) statt des wahren Werts. \textbf{Umformulierung:} Der Nenner vereinfacht sich analytisch zu $1 + \epsilon - 1 = \epsilon$, also
  \[
  f(\epsilon) = \frac{1}{\epsilon}.
  \]
  Diese Form enthält keine Subtraktion nahezu gleicher Zahlen mehr und liefert für $\epsilon = 10^{-20}$ korrekt
  \[
  f(10^{-20}) = \frac{1}{10^{-20}} = 10^{20}.
  \]
  \item Für $x = 10^{20}$ ist $x^2 = 10^{40} > 10^{38}$, also größer als die größte darstellbare Zahl in single precision: \textbf{Overflow} --- das Ergebnis wird als $+\infty$ approximiert. Für $y = 10^{-25}$ ist $y^2 = 10^{-50}$, eine Zahl so nah an null, dass sie in single precision auf $0$ gerundet wird: \textbf{Underflow}. Gefährliche Folgerechnung: Wird anschließend z.\,B.\ $y/y^2$ gebildet, dividiert man durch die gespeicherte $0$ statt durch $10^{-50}$ --- das Ergebnis wird $\infty$ (bzw.\ ein Laufzeitfehler) statt des wahren Werts $10^{-25}/10^{-50} = 10^{25}$, der in single precision problemlos darstellbar wäre ($10^{25} < 10^{38}$). Der Schaden entsteht hier also allein durch den Underflow im Nenner.
  \item Bei niedrigpräziser Arithmetik sind Machine Epsilon groß und die darstellbaren Wertemengen extrem klein (bei BNNs nur $\{-1, +1\}$). Rundungs- und Quantisierungsfehler sind dadurch um viele Größenordnungen größer als bei float/double, können sich in den vielen aufeinanderfolgenden Operationen eines neuronalen Netzes verstärken und akkumulieren und so das Inferenzergebnis spürbar verfälschen. Wer Modelle on the edge einsetzt, muss daher verstehen, welche Werte überhaupt darstellbar sind, wo gerundet wird und wie sich diese Fehler fortpflanzen --- sonst lassen sich Genauigkeitsverluste durch Quantisierung weder vorhersagen noch kontrollieren.
\end{enumerate}

\end{document}
